% rubber: setlist arguments --shell-escape

%\documentclass[handout]{beamer}
\documentclass[table]{beamer}
\usepackage{minted}
\usetheme{Luebeck}
\usepackage[utf8]{inputenc}

\newminted{cpp}{fontsize=\small}

\title[Visualizing Parallelism in CS 2]{Visualizing Parallelism in CS 2}
\author{Sean Massung \& Cinda Heeren}
\date{May 20th 2013}
\institute{University of Illinois at Urbana-Champaign}

\begin{document}
\frame{\titlepage}

\section{Introduction}

\begin{frame}
    \frametitle{A Spring 2012 Early Adopter}
    \begin{itemize}
         \item IEEE-TCPP Curriculum Initiative for CS 2 at UIUC
         \item Three main parallel programming concepts,
             each delivered during a two hour discussion/lab section
         \item Material is graded, and key topics are evaluated on exams
         \item Three lessons explore the basics of parallelism in a
             \textbf{visual} manner: using OpenMP, race conditions,
             and reductions
         \item Tethering images to understanding parallelism enhances
             student interest and provides an alternative route to
             mastery of the material
    \end{itemize}
\end{frame}

\section{New Labs}

\subsection{Lab 1: Intro to Parallelism}

\begin{frame}[fragile]
    \frametitle{Image manipulation}
    \begin{columns}
    \column{.45\textwidth}
    \begin{itemize}
        \item In the coding portion of the lab, we expose students
            to computation across threads
        \item For one task, students simply write a function that
            removes the green color component from an image
    \end{itemize}
    \column{.55\textwidth}
    \begin{cppcode}
#pragma omp parallel for
for(int i = 0; i < width; ++i)
{
   for(int j = 0; j < height; ++j)
   {
      *output(i,j) = *source(i,j);
       output(i,j)->green = 0;
   }
}
    \end{cppcode}
    \end{columns}
\end{frame}

\begin{frame}
    \frametitle{A visual partition of work}
    \begin{columns}
        \column{.5\textwidth}
        \includegraphics[width=2.2in]{semi-processed-remove.png}
        \column{.5\textwidth}
        \begin{itemize}
        \item Students are given augmented code that stops execution midway
            through the operation
        \item They see four threads operating on the
            image removing the green component
        \item The threads are partitioning the image along the width as
            shown in the code snippet
        \end{itemize}
    \end{columns}
\end{frame}

\subsection{Lab 2: Race Conditions}

\begin{frame}[fragile]
    \frametitle{Something's wrong with this image...}
    \begin{columns}
        \column{.5\textwidth}
        \begin{itemize}
            \item Students learn that correctly parallelizing
            programs does not just consist of blindly pasting a
            \verb|#pragma| on an outer for loop
            \item Image output from the provided code is shown to the right
        \end{itemize}
        \column{.5\textwidth}
        \includegraphics[width=2.2in]{flipped.png}
    \end{columns}
\end{frame}

\begin{frame}[fragile]
    \frametitle{The buggy code}
        \begin{cppcode}
RGBAPixel temp;
#pragma omp parallel for
for(int i = 0; i < width; ++i)
{
   for(int j = 0; j < height / 2; ++j)
   {
      temp = *image(i, j);
      *image(i, j) = *image(i, height - 1 - j);
      *image(i, height - 1 - j) = temp;
   }
}
        \end{cppcode}
        Moving the declaration of \verb|temp| inside the for loop
        fixes the problem, creating a local \verb|temp| variable for each
        thread.
\end{frame}

\subsection{Lab 3: Reductions}

\begin{frame}[fragile]
    \frametitle{We want to create a color histogram}
    \begin{columns}
        \column{.5\textwidth}
        \begin{itemize}
            \item The last lab section introduces a paradigm for solving complex
                data dependency issues, namely {\em reductions}
            \item We present reductions as a general algorithmic technique, so
                as to  provide a stepping stone to understanding the MapReduce
                programming model
        \end{itemize}
        \column{.5\textwidth}
        \includegraphics[width=2.2in]{chart.png}
    \end{columns}
\end{frame}

\begin{frame}[fragile]
    \frametitle{Creating the histogram}
    \begin{minted}[fontsize=\tiny]{cpp}
map<RGBAPixel, int> ret_freq;

#pragma omp parallel
{
   map<RGBAPixel, int> local_freq;
#pragma omp for
   for(int i = 0; i < width; ++i)
   {
      for(int j = 0; j < height; ++j)
          ++local_freq[ *image(i,j) ];
   }
#pragma omp critical
   {
      map<RGBAPixel, int>::iterator curr;
      for(curr = local_freq.begin(); curr != local_freq.end(); ++curr)
      {
         int count = curr->second;
         ret_freq[ curr->first ] += count;
      }
   }
}

return ret_freq;
    \end{minted}
\end{frame}

\section{Evaluation}

\begin{frame}
    \frametitle{Evaluating the students}
    \begin{itemize}
        \item Student performance on lab exercises was exemplary with average
            scores of 95\% on all three assignments
        \item Besides grading lab work, questions regarding parallelism were
            included on exams
        \item The following slides contain examples of the questions we have
            included on various tests over the past two semesters
    \end{itemize}
\end{frame}

\subsection{Question 1}

\begin{frame}[fragile]
    \frametitle{Pop quiz!}

    \begin{columns}

    \column{.5\textwidth}
    \begin{minted}[fontsize=\tiny]{cpp}
// (i)

#pragma omp parallel for
for (int i = 1; i < 100; i++)
    colorArray[i] = colorArray[i-1];

// (ii)

#pragma omp parallel for
for (int i = 0; i < width; i++)
{
    for (int j = 0; j < height/2; j++)
    {
        RGBAPixel temp = *img(i, j);
        *img(i,j) = *img(i, height-1-j);
        *img(i, height-1-j) = temp;
    }
}

// (iii)    

#pragma omp parallel for
for(int i = 0; i < 10; i++)
    for(int j = 0; j < 10; j++)
        table[i][j] = (i+1)*(j+1); 
    \end{minted}

    \column{.5\textwidth}
    Which of the code examples is/are NOT correctly parallelized?
    {\small
    \begin{enumerate}[(a)]
        \item Only item (i) is incorrect.
        \item Only item (ii) is incorrect.
        \item Only item (iii) is incorrect.
        \item Two of the above examples are incorrect.
        \item All statements (i), (ii), and (iii) are correct.
    \end{enumerate}
    }
    \end{columns}

\end{frame}

\subsection{Question 2}

\begin{frame}
    \frametitle{Reductions}
    What is the definition of the term \emph{reduction}?

    \begin{enumerate}[(a)]
        \item A reduction performs the same instructions on data
            across multiple threads.
        \item A reduction occurs when private data on individual
            threads is assembled into a general solution.
        \item Reduction is a technique wherein parallelism is
            applied to the portion of a program that requires
            the most computation. 
        \item Reduction is just another term for {\em speedup}.
        \item None of these is the correct choice.
    \end{enumerate}
\end{frame}

\subsection{Question 3}

\begin{frame}
    \frametitle{Amdahl's Law}
    Suppose an algorithm takes 7 seconds to run serially, and 2 seconds to run
    in parallel.  Then the {\em speedup} for the parallelized code is:

    \begin{enumerate}[(a)]
        \item $\frac{2}{7}$
        \item $\frac{7}{2}$
        \item $\frac{7-2}{7}$
        \item The speedup cannot be determined because the number of processors is not known.
        \item None of these answers is correct.
\end{enumerate}
\end{frame}

\section{Conclusions and Future Work}

\begin{frame}
    \frametitle{Conclusions}
    \begin{itemize}
    \item Foundations for including parallelism in introductory and second-level
        courses exist, but should be presented in a more approachable way
    \item Applying parallelism to images gives a useful purpose to students’
        work and immediate, visual feedback
    \item Watching threads work on images increases student understanding and
        allows them to connect with the assignment on a tangible level
    \end{itemize}
\end{frame}

\begin{frame}
    \frametitle{Future Work}
    \begin{itemize}
        \item Due to course reform, we will have more space added to the class:
            this time can be used to deepen the parallel concentration of the
            course in addition to examining existing topics in more depth
        \item A fourth lab section is under development that deals with OpenMP
            tasks and their use in parallel sorting algorithms
    \end{itemize}
\end{frame}


\end{document}
